% paper/sections/12h_error_budget.tex
%
% §12.6  精度バジェット（Error Budget）
%   §12.1--§12.5 の物理的整合性検証を踏まえた精度設計の総括
%   物理ベンチマーク（静止液滴・毛細管波・気泡上昇・RT不安定）は
%   第~\ref{sec:validation}章（§12）に収録．

% ═══════════════════════════════════════════════════════════════════════════════
\subsection{精度バジェット}
\label{sec:error_budget}
% ═══════════════════════════════════════════════════════════════════════════════

本章の検証結果を統合し，各コンポーネントの実測精度を一覧化する．
物理的な多相流ベンチマーク（毛細管波・気泡上昇・Taylor 液滴変形）は §~\ref{sec:validation} で扱う．

\begin{table}[htbp]
\centering
\caption{コンポーネント精度一覧（実測値）}
\label{tab:error_budget}
\renewcommand{\arraystretch}{1.35}
\begin{tabular}{>{\raggedright\arraybackslash}p{32mm}>{\raggedright\arraybackslash}p{22mm}>{\raggedright\arraybackslash}p{22mm}c>{\raggedright\arraybackslash}p{14mm}>{\raggedright\arraybackslash}p{32mm}}
\toprule
コンポーネント & 設計精度 & 実測精度 & 判定 & 検証箇所 & 誤差が影響する物理量 \\
\midrule
CCD 微分（内部） & $\Ord{h^6}$ & $\Ord{h^{5.97}}$ & \checkmark & §~\ref{sec:verify_ccd_convergence} & 速度微分，PPE 右辺 \\
CCD-PPE（変密度） & $\Ord{h^6}$ & $\Ord{h^{5.8}}$ & \checkmark & §~\ref{sec:verify_ppe_varrho} & 圧力場，速度補正 \\
CCD-PPE（$\rho_l/\rho_g{=}2$） & --- & $\Ord{h^{1\text{--}2}}$ & $\triangle$ & §~\ref{sec:varrho_grid_convergence} & 変密度圧力場 \\
HFE 2D 場延長 & $\Ord{h^6}$ & $\Ord{h^{5.8}}$\textsuperscript{$\P$} & $\triangle$ & §~\ref{sec:verify_field_extension} & IPC $\bnabla p^n$，速度補正 \\
DCCD CLS 移流 & $\Ord{\varepsilon_d h^2}$ & $\Ord{h^{2}}$\textsuperscript{$\dagger$} & \checkmark & §~\ref{sec:verify_cls_advection} & 界面位置，体積保存 \\
CCD 曲率 $\kappa$（CLS） & $\Ord{h^6}$\textsuperscript{$\|$} & $\Ord{h^{1.0}}$ & $\triangle$ & §~\ref{sec:verify_curvature} & 表面張力体積力 \\
AB2+IPC 時間（単相） & $\Ord{\Delta t^2}$ & $\Ord{\Delta t^{2.00}}$ & \checkmark & §~\ref{sec:tgv_temporal} & 全体時間精度 \\
NS 空間（Kovasznay） & $\Ord{h^4}$--$\Ord{h^6}$ & $\Ord{h^{3.97}}$ & \checkmark & §~\ref{sec:kovasznay} & NS 運動量残差 \\
CSF 表面張力 & $\Ord{\varepsilon^2}$ & $\Ord{h^{1.5\text{--}2}}$ & \checkmark & §~\ref{sec:bf_static_droplet} & Laplace 圧精度 \\
\midrule
\multicolumn{6}{l}{\textit{追加検証}} \\
\midrule
混合偏微分 $\partial_{xy}$ & $\Ord{h^6}$ & $\Ord{h^{6.0}}$（周期） & \checkmark & §~\ref{sec:verify_mixed_partial} & 交差粘性項 \\
CN 粘性項（対角） & $\Ord{\Delta t^2}$ & $\Ord{\Delta t^{2.0}}$ & \checkmark & §~\ref{sec:verify_cn_viscous} & 粘性時間積分 \\
交差粘性項（陽的） & $\Ord{\Delta t}$\textsuperscript{$\ddagger$} & $\Ord{\Delta t^{1.0}}$ & \checkmark & §~\ref{sec:verify_cross_viscous} & 界面近傍で律速 \\
CCD $D^{(2)}$ スペクトル半径 & --- & $\rho h^2 = 9.6$ & \checkmark & §~\ref{sec:verify_ccd_spectral_radius} & FD 比 2.4 倍 \\
DGR 厚み制御 & $\varepsilon_{\mathrm{eff}}/\varepsilon \approx 1$ & 1.03 & \checkmark & §~\ref{sec:verify_dgr} & 界面厚維持 \\
CCD vs FD 寄生流れ & CCD $\ll$ FD & FD/CCD $= 69\times$ & \checkmark & §~\ref{sec:force_balance} & Balanced-Force 効果 \\
交差粘性 CFL & $\Delta t \propto h^2/\Delta\mu$ & $C_{\mathrm{cross}} \approx 0.22$ & \checkmark & §~\ref{sec:cross_viscous_cfl} & 安定性制約 \\
毛細管 CFL & $\Delta t_\sigma \propto \Delta x^{3/2}$ & 指数 $= 1.500$，比率 $1.28$ & \checkmark & §~\ref{sec:advection_pressure_coupling} & 時間ステップ制限 \\
\bottomrule
\multicolumn{6}{l}{\footnotesize \textsuperscript{$\dagger$}DCCD 散逸パラメータ $\varepsilon_d$ によらず界面の CLS プロファイル幅が支配する誤差（§~\ref{sec:verify_cls_advection}）．} \\
\multicolumn{6}{l}{\footnotesize \textsuperscript{$\ddagger$}$\mu_l/\mu_g \gg 1$ で陽的交差微分項が $\Ord{\mu_{\max}\Delta t}$ の打切り誤差を生じる（§~\ref{sec:time_int}）．} \\
\multicolumn{6}{l}{\footnotesize \textsuperscript{$\P$}$N\le128$ で $p=6.33,\,5.79$；$N{=}256$ で補間精度床（$L_\infty{=}4.6\times10^{-11}$）に到達（条件付き合格；§~\ref{sec:verify_summary}）．} \\
\multicolumn{6}{l}{\footnotesize \textsuperscript{$\|$}CCD 演算子単体は $\Ord{h^6}$ だが，CLS パイプライン（$\varepsilon{=}1.5h$ 正則化＋$|\phi|<3h$ 評価帯）では $\Ord{h^1}$ が律速．} \\
\multicolumn{6}{l}{\footnotesize \checkmark: 設計精度達成，$\triangle$: 条件付き合格（表~\ref{tab:verification_summary} の判定基準に準拠）．} \\
\multicolumn{6}{l}{\footnotesize 検証箇所：§~\ref{sec:verify_ccd_convergence}--§~\ref{sec:verify_dgr} は第~\ref{sec:numerical_integrity}章（コンポーネント検証），それ以外は本章（物理的整合性検証）．} \\
\end{tabular}
\end{table}

% ─────────────────────────────────────────────────────────────────────────────
\paragraph{律速段階の特定}
\label{sec:accuracy_summary}

\begin{description}[style=nextline,leftmargin=2em]
  \item[\textbf{単相 NS ソルバ}（設計精度をほぼ完全に発現）]
    \begin{itemize}[noitemsep]
      \item 空間律速：CCD 境界スキーム $\Ord{h^5}$ $\times$ 非線形増幅 → 実効 $\Ord{h^{3.97}}$（§~\ref{sec:kovasznay}）
      \item 時間律速：AB2+IPC $\Ord{\Delta t^2}$
    \end{itemize}

  \item[\textbf{二相流（$\rho_l/\rho_g \leq 5$，smoothed Heaviside）}]
    \begin{itemize}[noitemsep]
      \item 空間律速：\textbf{CSF モデル誤差 $\Ord{h^2}$}（CCD-PPE $\Ord{h^6}$ は過剰）
      \item 変密度一括 PPE：界面近傍で $\Ord{h^{1\text{--}2}}$ に劣化（CSF と同次）
      \item 静止界面：200 ステップ安定（表~\ref{tab:static_longterm}）
      \item 移動界面＋$\sigma > 0$：\textbf{IPC 発散}（$\bnabla p^n$ が界面ジャンプを横断）
    \end{itemize}

  \item[\textbf{PPE 離散化}]
    \begin{itemize}[noitemsep]
      \item CCD PPE 欠陥補正は $L_L^{-1} L_H$（$L_L$: $\Ord{h^2}$ FD 基底ソルバ，$L_H$: CCD 高次演算子）の不定性で収束不能 → FD 直接法 + CCD 勾配を採用
      \item 投影不整合 $\Ord{h^2}$：CSF と同次のため全体精度を律速しない
    \end{itemize}
\end{description}

\noindent\textbf{二相流（高密度比・界面型密度場）：}
変密度一括 PPE の欠陥補正は再試験では $\rho_l/\rho_g = 100$ まで発散しなかったが，
高解像度側で漸近格子収束が失われる（§~\ref{sec:varrho_grid_convergence}）．
分相 PPE + DC $k{=}3$ は各相内部で $\Ord{h^{7.0}}$ を回復する
（§~\ref{sec:verify_split_ppe_dc}）．

% ─────────────────────────────────────────────────────────────────────────────
\paragraph{CCD 高次精度の意義}

CCD の $\Ord{h^6}$ 精度は現状の CSF 近似（$\Ord{h^2}$）のもとでは
\textbf{直接的な全体精度向上には寄与しない}．
しかし以下の 2 つの経路で本質的な役割を果たす：
\begin{enumerate}[noitemsep]
  \item \textbf{寄生流れの係数低減：}
        Balanced-Force 条件下では収束\textbf{次数}は CSF 律速で $\Ord{h^2}$ だが，
        収束\textbf{係数}は曲率精度に比例する．
        CCD（$p_\kappa = 6$）は 2 次法と比べ $h = 1/64$ で
        $(1/64)^4 \approx 10^{-7}$ 倍の係数低減効果を持つ
        （§~\ref{sec:balanced_force}）．
  \item \textbf{分相 PPE 解法による精度回復：}
        分相 PPE 解法は CSF の $\Ord{h^2}$ モデル誤差を原理的に除去し，
        各相の PPE を独立に解くことで
        CCD の $\Ord{h^6}$ 設計精度を回復する
        （§~\ref{sec:interface_crossing}）．
\end{enumerate}

% ─────────────────────────────────────────────────────────────────────────────
\paragraph{ソルバの適用範囲}

表~\ref{tab:usage_constraints} に，本章の検証結果から導かれるソルバの適用制約を集約する．

\begin{table}[htbp]
\centering
\caption{現行ソルバの適用制約（本章の検証に基づく）}
\label{tab:usage_constraints}
\renewcommand{\arraystretch}{1.3}
\small
\begin{tabular}{p{42mm}ccp{38mm}}
\toprule
条件 & 対応 & 参照 & 備考 \\
\midrule
単相 NS（任意 Re） & \checkmark & §~\ref{sec:ns_time_accuracy} & $\Ord{\Delta t^2}$，$\Ord{h^4}$ 実効精度 \\
静止界面＋$\sigma > 0$，$\rho_l/\rho_g \leq 5$ & \checkmark & §~\ref{sec:twophase_static_longterm} & 200 ステップ安定 \\
移動界面，$\sigma = 0$，$\rho_l/\rho_g \leq 5$ & \checkmark & §~\ref{sec:galilean_invariance} & Galilean 不変，RT 誤差 2.8\% \\
移動界面＋$\sigma > 0$（任意 $\rho$） & $\times$ & §~\ref{sec:coupling_limitations} & IPC $\bnabla p^n$ が界面跳びで発散 \\
高密度比 smoothed H 一括 PPE & $\triangle$ & §~\ref{sec:interface_crossing} & DC 発散は再現せず；格子収束はプラトー化 \\
分相 PPE + DC $k{=}3$ & \checkmark & §~\ref{sec:verify_split_ppe_dc} & 全密度比で $\Ord{h^{7.0}}$ \\
\bottomrule
\end{tabular}
\end{table}

% ─────────────────────────────────────────────────────────────────────────────
\paragraph{§~\ref{sec:validation}章への接続}

§~\ref{sec:twophase_static_longterm} で低密度比の静止界面二相流が安定動作することを確認し，
§~\ref{sec:galilean_invariance} で移動界面＋表面張力には圧力場の分相延長が必要であることを
特定した．
§~\ref{sec:interface_crossing} で smoothed Heaviside 一括解法の格子収束劣化を定量化し，
§~\ref{sec:verify_split_ppe_dc} では分相 PPE + DC $k{=}3$ により
$\rho_l/\rho_g = 1000$ でも $\Ord{h^{7.0}}$ 格子収束が達成されることを実証した
（表~\ref{tab:split_ppe_dc_convergence}）．
第~\ref{sec:validation}章では本章で確立した動作範囲内
（$\rho_l/\rho_g \leq 10$，$\sigma$ 有限で静的または低速移動界面）において
物理ベンチマーク 3 件を実施する：
毛管波の振動減衰（Prosperetti 解析解との比較），
気泡上昇（Hysing ら 2009 Test Case 1），
および線形せん断流中の液滴変形（Taylor 1932）．
いずれも表面張力・粘性・浮力の結合を定量検証するものであり，
ただし §~\ref{sec:val_grid_policy} で述べる通り，
非一様格子（$\alpha > 1$）は計量変換誤差により不安定であるため，
全実験を一様格子（$\alpha = 1$）で実施する．

\noindent\textbf{本章の残存限界（§~\ref{sec:validation}章で補完予定）：}
\begin{itemize}[noitemsep]
  \item \textbf{移動界面＋$\sigma > 0$ の結合安定性：}
        全 $\sigma > 0$ テストは静止界面，全移動界面テストは $\sigma = 0$．
        両者の結合（界面変形を伴う表面張力流れ）は
        §~\ref{sec:validation} の毛管波・気泡上昇で初めて検証される．
  \item \textbf{分相 PPE の物理ベンチマーク：}
        §~\ref{sec:verify_split_ppe_dc} は製造解（MMS）のみ．
        分相 PPE の NS パイプライン組み込みは §~\ref{sec:validation} で実施する．
  \item \textbf{再初期化--移流カップリング：}
        §~\ref{sec:verify_reinit_shift} の再初期化シフト計測は初期場のみ．
        200 ステップ以上の長時間 NS 解での蓄積効果は
        §~\ref{sec:validation} の体積保存計測で検出される．
  \item \textbf{$\mu_l/\mu_g \gg 1$ のシステムレベル時間精度：}
        §~\ref{sec:verify_cross_viscous} で成分レベル $\Ord{\Delta t}$ を確認済み．
        Taylor 変形テスト（§~\ref{sec:validation}）で高粘性比の結合効果を検証予定．
  \item \textbf{毛細管 CFL 安定限界の直接検証：}
        §~\ref{sec:advection_pressure_coupling} で CFL 指数を計測済みだが，
        限界値を超えた際の発散を示す安定境界テストは未実施．
        §~\ref{sec:validation} の毛管波テストで間接検証される．
\end{itemize}

\noindent\textbf{検証対象外の保存量：}
角運動量はコロケート格子上の離散 NS ソルバでは一般に厳密保存されない%
\footnote{スタガード格子は運動量の局所保存を格子構造により保証するが，
コロケート格子では圧力--速度の補間構造上，
角運動量フラックスの厳密な保存は保証されない．
本稿では DCCD フィルタ（$\varepsilon_d = 1/4$）による欠陥補正で
チェッカーボード不安定性を抑制しており（§~\ref{sec:dccd_decoupling}），
面流速補間に起因する角運動量誤差は副次的である．
角運動量保存の厳密な検証にはスタガード格子への拡張が必要であり，
本稿の範囲外とする．}．
したがって本章では角運動量保存の独立検証は実施しない．
